\chapter{Analizador léxico}

El analizador léxico transforma el archivo fuente en una secuencia de tokens.
Cada token conserva su tipo, lexema y posición de inicio, de modo que las etapas
posteriores puedan reconocer las construcciones del lenguaje e informar errores.
Los espacios en blanco y los comentarios se descartan durante este proceso.

% ############################################################
% #                                                          #
% #               ALFABETO DE ENTRADA                        #
% #                                                          #
% ############################################################

\section{Alfabeto de entrada}
Se denomina $\Sigma$ al conjunto de caracteres considerado para describir las
reglas léxicas. La pertenencia al alfabeto no implica que un carácter pueda
aparecer en cualquier contexto: por ejemplo, \verb~@~ puede aparecer dentro de
un comentario o una cadena, pero no constituye un token válido por sí solo.
\begin{itemize}[leftmargin=*]
  \item Letras minúsculas: abcdefghijklmnñopqrstuvwxyz
  \item Letras mayúsculas: ABCDEFGHIJKLMNÑOPQRSTUVWXYZ
  \item Números: 0123456789
  \item Símbolos: \verb~+-*/=<>!&|;:,.$#%@^¿?¡(){}[]_~
  \item Comillas y barra invertida: \verb~"~ (delimitador de cadenas),
        \verb~'~ (comilla simple) y \verb~\~ (inicio de una secuencia de escape).
  \item Escapes: espacio, \verb~\n~, \verb~\r~, \verb~\t~ y \verb~\v~.
  \item En los literales \texttt{Str} se admiten además las vocales acentuadas:
        áéíóúÁÉÍÓÚ.
\end{itemize}
Las restricciones particulares de las cadenas se detallan en su definición.
El fin de archivo (EOF) es una condición de lectura, no un carácter de $\Sigma$.

% ############################################################
% #                                                          #
% #               TOKENS RECONOCIDOS                         #
% #                                                          #
% ############################################################

\section{Tokens reconocidos}

% ============================================================
%                IDENTIFICADOR MET_VAR
% ============================================================

\subsection{Identificador de método o variable}
\begin{itemize}[leftmargin=*]
  \item Expresión regular: \verb~[a-z] . ([a-z] | [A-Z] | [0-9] | [_])*~
  \item Tipo de token: \lexcode{IDENTIFIER_VAR_METHOD}.
  \item Ejemplos: \lexcode{nombre}, \lexcode{numero_1}.
\end{itemize}
Se consume el identificador completo y luego se consulta la tabla de palabras
reservadas. Por ello, \lexcode{if} es una palabra reservada, mientras que
\lexcode{ifElse} es un identificador.

% ============================================================
%                IDENTIFICADOR CLASE
% ============================================================

\subsection{Identificador de clase}
\begin{itemize}[leftmargin=*]
  \item Expresión regular: \verb~[A-Z] . ([a-z] | [A-Z] | [0-9] | [_])*~
  \item Tipo general: \lexcode{IDENTIFIER_CLASS}.
  \item Ejemplos: \lexcode{Persona}, \lexcode{Clase_1}.
\end{itemize}
Los nombres \lexcode{Int}, \lexcode{Str}, \lexcode{Bool} y \lexcode{Array}
se reclasifican mediante la tabla de tipos predefinidos.

% ============================================================
%                PALABRAS RESERVADAS
% ============================================================

\subsection{Palabras reservadas}
Salvo las constantes booleanas, cada palabra se asocia al tipo
\lexcode{KW_} seguido de su nombre en mayúsculas; por ejemplo,
\lexcode{class} produce \lexcode{KW_CLASS} y \lexcode{void} produce
\lexcode{KW_VOID}. Los lexemas \lexcode{true} y \lexcode{false} producen
\lexcode{BOOLEAN_CONSTANT}. El literal nulo \lexcode{nil} produce
\lexcode{KW_NIL}.
\begin{itemize}[leftmargin=*]
  \item Expresiones regulares:
  \subitem \verb~class: c.l.a.s.s~
  \subitem \verb~if: i.f~
  \subitem \verb~impl: i.m.p.l~
  \subitem \verb~in: i.n~
  \subitem \verb~else: e.l.s.e~
  \subitem \verb~false: f.a.l.s.e~
  \subitem \verb~fn: f.n~
  \subitem \verb~for: f.o.r~
  \subitem \verb~ret: r.e.t~
  \subitem \verb~while: w.h.i.l.e~
  \subitem \verb~true: t.r.u.e~
  \subitem \verb~new: n.e.w~
  \subitem \verb~nil: n.i.l~
  \subitem \verb~st: s.t~
  \subitem \verb~self: s.e.l.f~
  \subitem \verb~pub: p.u.b~
  \subitem \verb~void: v.o.i.d~
\end{itemize}

% ============================================================
%                LITERALES ENTEROS
% ============================================================

\subsection{Literales enteros}
\begin{itemize}[leftmargin=*]
  \item Expresión regular: \verb~[0-9]+~
  \item Tipo de token: \lexcode{INT_CONSTANT}.
\end{itemize}
El signo no forma parte del literal: \verb~-42~ se reconoce como el operador
de resta seguido del entero \verb~42~. En esta etapa se conserva el lexema
sin convertirlo a un valor numérico ni comprobar su rango.

% ============================================================
%                LITERALES STRING
% ============================================================

\subsection{Literales string}
\shorthandoff{"}
Sea $C$ el conjunto de caracteres ordinarios admitidos dentro de una cadena:
letras, dígitos, símbolos del alfabeto, espacio, tabulador y vocales acentuadas,
excluyendo la comilla doble, la barra invertida, los saltos de línea y el
carácter nulo. Las comillas dobles y las barras invertidas internas se
representan mediante escapes. La expresión regular es:
\[
\texttt{"}\cdot
\left(C\mid
\backslash\cdot
(\texttt{n}\mid\texttt{t}\mid\texttt{r}\mid
\texttt{"}\mid\texttt{'}\mid\backslash)
\right)^*\cdot\texttt{"}
\]
\begin{itemize}[leftmargin=*]
  \item Tipo de token: \lexcode{STRING_CONSTANT}.
  \item Se admiten los escapes \verb~\n~, \verb~\t~, \verb~\r~,
        \verb~\"~, \verb~\'~ y \verb~\\~.
  \item La cadena puede estar vacía y debe cerrarse en la misma línea física.
        El escape \verb~\n~ permite representar un salto de línea en su valor.
  \item El tamaño máximo es de 1024 caracteres de contenido.
  \item No se admite el carácter nulo ni llegar al fin de archivo sin cerrar
        la cadena.
\end{itemize}
El token conserva las comillas delimitadoras y los escapes tal como aparecen
en el archivo. Actualmente, el límite se controla sobre el lexema sin sus
comillas, por lo que cada escape ocupa dos posiciones en el contador.
\shorthandon{"}

% ============================================================
%                ESPACIOS EN BLANCO
% ============================================================

\subsection{Espacios en blanco}
\begin{itemize}[leftmargin=*]
  \item Expresión regular para una secuencia no vacía:
        \verb~(" " | \n | \t | \r | \v)+~
\end{itemize}
Estos caracteres separan lexemas y no producen tokens. La implementación lee
líneas mediante \lexcode{BufferedReader.readLine()}, que elimina sus
terminadores, y descarta los restantes espacios con
\lexcode{Character.isWhitespace()}.

% ============================================================
%                COMENTARIOS
% ============================================================

\subsection{Comentarios}
Sea $U$ el conjunto de caracteres que pueden aparecer en el contenido de un
comentario. La implementación los omite sin aplicar el filtro de
\lexcode{Alphabet}. Definimos $A=U-\{\texttt{*}\}$ y $B=U-\{\texttt{*},\texttt{/}\}$.
\begin{itemize}[leftmargin=*]
  \item Expresión regular para comentario simple:
  \[
  \texttt{//}\cdot U^*
  \]
  Termina al finalizar la línea o el archivo; no necesita un salto de línea
  final para ser válido.
  \item Expresión regular para comentario multilínea:
  \[
  \texttt{/*}\cdot
  (A\mid\texttt{*}^{+}\cdot B)^*\cdot
  \texttt{*}^{+}\cdot\texttt{/}
  \]
  Termina en el primer cierre \verb~*/~. No se admite anidamiento y llegar
  al fin del archivo antes del cierre genera un error léxico.
\end{itemize}
Los comentarios no producen tokens. Dentro de una cadena, los caracteres
\verb~//~ y \verb~/*~ forman parte del contenido y no inician comentarios.

% ============================================================
%                CLASES PREDEFINIDAS DE TINYS
% ============================================================

\subsection{Clases predefinidas}
Debe distinguirse entre el nombre de una clase predefinida y su clasificación
léxica. En la implementación, \lexcode{Int}, \lexcode{Str}, \lexcode{Bool}
y \lexcode{Array} producen, respectivamente, \lexcode{INT_CLASS},
\lexcode{STR_CLASS}, \lexcode{BOOL_CLASS} y \lexcode{ARRAY_CLASS}.
Los nombres \lexcode{Object}, \lexcode{IO} e \lexcode{Iterator} producen
\lexcode{IDENTIFIER_CLASS}; el tratamiento de las entidades predefinidas
corresponde a las etapas posteriores.
\begin{itemize}[leftmargin=*]
  \item Expresiones regulares:
  \subitem \verb~Object: O.b.j.e.c.t~
  \subitem \verb~IO: I.O~
  \subitem \verb~Iterator: I.t.e.r.a.t.o.r~
  \subitem \verb~Array: A.r.r.a.y~
  \subitem \verb~Int: I.n.t~
  \subitem \verb~Str: S.t.r~
  \subitem \verb~Bool: B.o.o.l~
\end{itemize}

% ============================================================
%                OPERADORES DE TINYS
% ============================================================

\subsection{Operadores}
\begin{itemize}[leftmargin=*]
  \item Expresiones regulares:

  \subitem \textbf{Operadores aritméticos:}
  \subitem \verb~Suma: +~
  \subitem \verb~Resta: -~
  \subitem \verb~Multiplicación: *~
  \subitem \verb~División: /~
  \subitem \verb~Incremento: +.+~
  \subitem \verb~Decremento: -.-~

  \subitem \textbf{Operadores relacionales:}
  \subitem \verb~Menor: <~
  \subitem \verb~Menor o igual: <.=~
  \subitem \verb~Mayor: >~
  \subitem \verb~Mayor o igual: >.=~
  \subitem \verb~Igualdad: =.=~
  \subitem \verb~Desigualdad: !.=~

  \subitem \textbf{Operadores lógicos:}
  \subitem \verb~AND: &.&~
  \subitem \verb~OR: |.|~
  \subitem \verb~NOT: !~

  \subitem \textbf{Operador de asignación:}
  \subitem \verb~Asignación: =~
\end{itemize}

% ============================================================
%                SÍMBOLOS DE TINYS
% ============================================================

\subsection{Símbolos}
\begin{itemize}[leftmargin=*]
  \item Expresiones regulares:
  \subitem \verb~Paréntesis izquierdo: (~
  \subitem \verb~Paréntesis derecho: )~
  \subitem \verb~Llave izquierda: {~
  \subitem \verb~Llave derecha: }~
  \subitem \verb~Corchete izquierdo: [~
  \subitem \verb~Corchete derecho: ]~
  \subitem \verb~Punto y coma: ;~
  \subitem \verb~Coma: ,~
  \subitem \verb~Punto: .~
  \subitem \verb~Dos puntos: :~
\end{itemize}

\subsection{Fin de archivo}
Al terminar el archivo, se devuelve un token de tipo \lexcode{EOF} y lexema
\lexcode{EOF}. Este token no tiene una expresión regular sobre el alfabeto:
se genera al detectar que no quedan líneas por leer. La palabra escrita
\lexcode{EOF} en un archivo fuente se clasifica como identificador de clase.
El valor \lexcode{UNDEFINED} se utiliza internamente para símbolos rechazados;
no se entrega como un token válido al analizador sintáctico.

\section{Diseño e implementación}
La implementación se divide en dos módulos: el analizador, ubicado en el
paquete \lexcode{uncuyo.compiladores.lexer}, y el ejecutador
\lexcode{uncuyo.compiladores.Etapa1}. El primero ofrece tokens a demanda;
el segundo reúne la secuencia y presenta el resultado por pantalla o archivo.
Esta separación permite reutilizar el mismo analizador en las etapas posteriores.

\subsection{Clases y responsabilidades}
\begin{description}[style=nextline]
  \item[\lexcode{Lexer}]
  Abre el archivo con un \lexcode{BufferedReader}, conserva la línea y posición
  actuales, omite espacios y solicita la clasificación del siguiente lexema.
  \lexcode{getNextToken()} devuelve un token por llamada y
  \lexcode{getAllTokens()} reúne la secuencia hasta incluir EOF.
  \item[\lexcode{TokenClassifier}]
  Implementa mediante métodos y ciclos las transiciones necesarias para
  reconocer identificadores, enteros, cadenas, operadores, símbolos y
  comentarios. Comparte el estado de lectura del \lexcode{Lexer}.
  \item[\lexcode{TokenDefinitions}]
  Mantiene mapas entre lexemas y tipos de token para palabras reservadas,
  operadores, símbolos, tipos predefinidos y símbolos no definidos.
  \item[\lexcode{Alphabet}]
  Agrupa caracteres y ofrece comprobaciones de pertenencia y de vocales
  acentuadas. El reconocimiento de cadenas utiliza estas comprobaciones.
  \item[\lexcode{Token} y \lexcode{TokenType}]
  El primero almacena lexema, tipo, fila y columna de inicio. El segundo es
  la enumeración de categorías que comparten las etapas del compilador.
  \item[\lexcode{Etapa1}]
  Interpreta los argumentos, solicita tokens hasta EOF y escribe el listado.
  \item[Excepciones]
  \lexcode{LexicalException} comunica errores léxicos.\par
  \lexcode{LexerFileReadException} comunica problemas de lectura del archivo.
\end{description}

\subsection{Posiciones y representación de lexemas}
La primera línea del archivo se enumera con uno. El índice interno de columna
comienza en cero. Los lexemas se conservan como cadenas de caracteres, incluyendo los delimitadores y escapes de los
literales \texttt{Str}.

\section{Manejo de errores léxicos}
La política es detener el análisis ante el primer error mediante una excepción
de tipo \lexcode{LexicalException}. No se intenta recuperar la lectura ni producir
tokens posteriores al error. El ejecutador reúne todos los tokens antes de emitir el
listado de éxito, por lo que un error impide mostrar un listado parcial exitoso.

\subsection{Errores detectados}
\begin{itemize}[leftmargin=*]
  \item Símbolos registrados como no definidos, como \verb~@~, \verb~$~,
        \verb~#~, \verb~%~, \verb~?~, \verb~¿~, \verb~^~ y \verb~_~ al inicio
        de un lexema; también \verb~&~ y \verb~|~ aislados.
  \item Vocales acentuadas dentro del recorrido de identificadores que comienza
        con una letra minúscula.
  \item Comentario multilínea sin cierre al finalizar el archivo.
  \item Cadena sin cerrar antes del final de su línea.
  \item Escape incompleto o no admitido, como \verb~\x~.
  \item Carácter nulo o carácter fuera del conjunto permitido dentro de una cadena.
  \item Cadena que supera el límite de longitud controlado por el reconocedor.
\end{itemize}
Los errores de apertura y lectura son problemas de entrada y salida, y se
distinguen de un programa con contenido léxicamente incorrecto.

\section{Salida del analizador léxico}
El procedimiento de compilación y ejecución se describe en la
sección~\ref{sec:compilacion-ejecucion} de la introducción.
Cada fila contiene el tipo, el lexema y la posición del token. El listado
incluye EOF. Al imprimir, \lexcode{TokenType.toString()} muestra los nombres
\texttt{Int}, \texttt{Str} y \texttt{Bool} para las categorías de constantes.
El encabezado de éxito es \texttt{CORRECTO: ANALISIS LEXICO}.

\section{Escenarios y casos de prueba}
Los archivos se encuentran en la carpeta \lexcode{tc/etapa1}, dentro de
la raíz del proyecto, y se distribuyen en dos directorios:
\begin{itemize}[leftmargin=*]
  \item Casos válidos: \lexcode{tc/etapa1/success}, con 17 archivos de prueba.
  \item Casos inválidos: \lexcode{tc/etapa1/error}, con 10 archivos de prueba.
\end{itemize}
Los escenarios cubren familias de tokens, separadores, comentarios,
cadenas con escapes, límites de longitud y entradas rechazadas.
Esta organización utiliza el directorio \lexcode{tc} solicitado
por el enunciado.

\subsection{Técnica y alcance de la verificación}
Las pruebas de esta etapa se implementan en \lexcode{LexerTest}.
Los casos se agrupan en particiones
de entradas aceptadas y rechazadas. En los casos válidos, \lexcode{assertDoesNotThrow} comprueba que se
alcance el final del análisis; en los inválidos, \lexcode{assertThrows} exige
una \lexcode{LexicalException}. Desde la raíz del proyecto, se ejecutan mediante:
\begin{verbatim}
mvn -Dtest=LexerTest test
\end{verbatim}

\subsection{Casos que deben aceptarse}
En cada caso se espera que el análisis llegue a EOF sin excepciones.
\begingroup
\small
\renewcommand{\arraystretch}{1.2}
\setlength{\tabcolsep}{5pt}
\setlength{\LTleft}{0pt}
\setlength{\LTright}{0pt plus 1fill}
\begin{longtable}{|
>{\raggedright\arraybackslash}p{0.40\textwidth}|
>{\raggedright\arraybackslash}p{0.28\textwidth}|
>{\raggedright\arraybackslash}p{\dimexpr0.32\textwidth-6\tabcolsep-4\arrayrulewidth\relax}|}
\caption{Casos válidos del analizador léxico}\label{tab:lexico-validos}\\
\hline
\textbf{Archivo} & \textbf{Propósito} & \textbf{Resultado esperado} \\
\hline
\endfirsthead
\hline
\textbf{Archivo} & \textbf{Propósito} & \textbf{Resultado esperado} \\
\hline
\endhead
\lexcode{adjacentTokens.s} & Reconocer lexemas consecutivos sin espacios. & Tokens correctamente separados y EOF. \\ \hline
\lexcode{colon.s} & Reconocer el carácter de dos puntos. & \lexcode{SPECIAL_SYMBOL_CL} y EOF. \\ \hline
\lexcode{comments.s} & Procesar comentarios simples y multilínea. & Sólo el entero, la cadena y EOF. \\ \hline
\lexcode{escapedSingleQuoteString.s} & Validar el escape de comilla simple. & Una cadena válida y EOF. \\ \hline
\lexcode{escapeStrings.s} & Validar escapes de salto, tabulación, comillas, barra y retorno. & Cadenas válidas y EOF. \\ \hline
\lexcode{identifiers.s} & Reconocer nombres con letras, dígitos y guion bajo. & Identificadores de clase o variable y método, y EOF. \\ \hline
\lexcode{identifiersAndCase.s} & Comprobar la distinción entre mayúsculas y minúsculas. & Clasificación correcta de cada lexema. \\ \hline
\lexcode{integers.s} & Reconocer cero y secuencias de varios dígitos. & Constantes enteras y EOF. \\ \hline
\lexcode{keywordPrefixes.s} & Reconocer nombres que comienzan con una palabra reservada. & Identificadores completos, sin dividir sus prefijos. \\ \hline
\lexcode{keywords.s} & Reconocer palabras reservadas y valores booleanos. & Palabras reservadas y booleanos bien clasificados. \\ \hline
\lexcode{operators.s} & Reconocer operadores de uno y dos caracteres. & Operadores correctamente agrupados y EOF. \\ \hline
\lexcode{specialSymbols.s} & Reconocer delimitadores y signos de puntuación. & Tokens de símbolos especiales y EOF. \\ \hline
\lexcode{stringContent.s} & Validar marcas de comentario, tildes, eñes y escapes en cadenas. & El contenido se conserva dentro de los literales. \\ \hline
\lexcode{stringLengthLimit.s} & Comprobar el máximo permitido de longitud. & Una cadena de 1024 caracteres aceptada y EOF. \\ \hline
\lexcode{strings.s} & Reconocer cadenas vacías, consecutivas y con distintos contenidos. & Constantes de cadena válidas y EOF. \\ \hline
\lexcode{unescapedSingleQuoteString.s} & Validar la comilla simple sin escape. & Una cadena válida y EOF. \\ \hline
\lexcode{whitespaceAndComments.s} & Descartar separadores y comentarios con símbolos. & Sólo el entero final y EOF. \\ \hline
\end{longtable}
\endgroup

\subsection{Casos que deben rechazarse}
En cada caso se espera una \lexcode{LexicalException} por el motivo indicado.
\begingroup
\small
\renewcommand{\arraystretch}{1.2}
\setlength{\tabcolsep}{5pt}
\setlength{\LTleft}{0pt}
\setlength{\LTright}{0pt plus 1fill}
\begin{longtable}{|
>{\raggedright\arraybackslash}p{0.40\textwidth}|
>{\raggedright\arraybackslash}p{0.28\textwidth}|
>{\raggedright\arraybackslash}p{\dimexpr0.32\textwidth-6\tabcolsep-4\arrayrulewidth\relax}|}
\caption{Casos inválidos del analizador léxico}\label{tab:lexico-invalidos}\\
\hline
\textbf{Archivo} & \textbf{Propósito} & \textbf{Resultado esperado} \\
\hline
\endfirsthead
\hline
\textbf{Archivo} & \textbf{Propósito} & \textbf{Resultado esperado} \\
\hline
\endhead
\lexcode{dollarSymbol.s} & Validar el símbolo de dólar fuera de cadenas y comentarios. & Error por símbolo no reconocido. \\ \hline
\lexcode{error.s} & Detectar la ausencia de la comilla de cierre. & Error por cadena sin cerrar. \\ \hline
\lexcode{invalidEscape.s} & Validar la secuencia \texttt{\textbackslash x}. & Error por escape no admitido. \\ \hline
\lexcode{leadingUnderscore.s} & Comprobar el inicio de un nombre con guion bajo. & Error por símbolo inicial no reconocido. \\ \hline
\lexcode{longitudStrInvalida.s} & Comprobar una cadena que excede el máximo. & Error por superar los 1024 caracteres. \\ \hline
\lexcode{singleAmpersand.s} & Validar un ampersand aislado. & Error por símbolo no reconocido. \\ \hline
\lexcode{singlePipe.s} & Validar una barra vertical aislada. & Error por símbolo no reconocido. \\ \hline
\lexcode{unclosedMultilineComment.s} & Detectar EOF antes de cerrar un comentario multilínea. & Error por comentario sin cerrar. \\ \hline
\lexcode{unclosedString.s} & Detectar el fin de línea dentro de una cadena sin cerrar. & Error por cadena sin cerrar. \\ \hline
\lexcode{undefinedSymbol.s} & Validar la arroba fuera de cadenas y comentarios. & Error por símbolo no reconocido. \\ \hline
\end{longtable}
\endgroup

